% IEEE RA-L full development draft, revised 2026-07-16.
% Magenta/red placeholders require measured evidence before submission.
\documentclass[letterpaper,10pt,conference]{ieeeconf}
\IEEEoverridecommandlockouts
\overrideIEEEmargins

\usepackage{amsmath,amssymb}
\usepackage{graphicx}
\usepackage{booktabs}
\usepackage{multirow}
\usepackage{xcolor}
\usepackage{cite}
\usepackage{url}

\newcommand{\TODO}[1]{\textcolor{red}{\textbf{[TODO: #1]}}}
\newcommand{\RESULT}[1]{\textcolor{magenta}{\textbf{[RESULT: #1]}}}
\newcommand{\obs}{\mathbf{o}}
\newcommand{\act}{\mathbf{a}}
\newcommand{\goal}{g}
\newcommand{\hid}{\mathbf{h}}
\newcommand{\qcmd}{\mathbf{q}^{\mathrm{cmd}}}
\newcommand{\qnom}{\mathbf{q}^{\mathrm{nom}}}
\newcommand{\qmeas}{\mathbf{q}^{\mathrm{meas}}}

\title{Progress-Aware Multimodal Autonomy for Robotic Bronchoscopy:\\
Recurrent Imitation Learning, Visual Junction Guidance, and Continuous Suction}

\author{Anonymous Authors%
\thanks{\TODO{Insert authors, affiliations, funding, and corresponding-author details.}}}

\begin{document}
\maketitle
\thispagestyle{empty}
\pagestyle{empty}

\begin{abstract}
Autonomous bronchoscopy requires long-horizon motor coordination, reliable
anatomical localization, local centering at visually ambiguous bifurcations,
and uninterrupted secretion management. We present a three-layer multimodal
autonomy framework learned from 87 real-robot expert demonstrations containing
30,730 synchronized frames. A goal-conditioned gated recurrent unit (GRU)
encodes an eight-frame history of 20-dimensional motor, perception, and depth
features and predicts three-axis motor increments together with an auxiliary
task label. The best saved model contains 773,221 parameters and achieved a
held-out action mean-squared error of 0.002084 and task accuracy of 100\%.
Layer~2 couples lumen geometry with a continuous progress--distance authority
schedule, allowing bounded visual steering to increase gradually without
overriding the global navigation direction. Layer~3 combines U-Net mucus
segmentation, YOLO-based anatomical recognition, Depth Anything V2 distance
estimation, continuously active vacuum suction, and progress-aware mission
continuity. Integrated phantom trials completed target-directed navigation,
online junction centering, secretion clearance, continuation of the remaining
inspection, and autonomous homing. Quantitative route, collision, clearance,
and timing results are organized in a locked reporting protocol and remain to
be inserted from the final trial sheet. The principal innovation is a unified
architecture in which learned global intent, geometry-driven local correction,
anatomical state, and therapeutic interaction operate at explicitly bounded
levels of authority.
\end{abstract}

\begin{keywords}
Medical robots and systems; imitation learning; autonomous navigation;
bronchoscopy; visual servoing; semantic segmentation.
\end{keywords}

\section{Introduction}

Robotic bronchoscopy improves stability and peripheral reach, with feasibility,
multicenter, and shape-sensing studies demonstrating its clinical potential
\cite{fielding2019firsthuman,chaddha2019multicenter,rojassolano2018feasibility,
chen2021benefit,kalchiem2022shapesensing}. Meta-analyses nevertheless report
heterogeneous performance and motivate standardized system-level evaluation
\cite{ali2023metaanalysis,zhang2024metaanalysis}. Most deployed platforms
remain operator driven, while autonomous traversal must handle repeated visual
patterns, uncertain lumen geometry, nonrigid contact, and secretions.

Recent bronchoscopy systems use simulator-trained shared control, multimodal
reinforcement learning, topological localization, or pure-vision agents
\cite{zhang2024aicopilot,zhao2024bronchocopilot,cold2026bronchobot,
wu2026longshort,tomasini2025topological}. Their progress exposes a central
control dilemma: a global policy is needed to preserve anatomical intent, but
local image feedback is needed to reduce wall contact. Unconstrained image
centering can select the wrong lumen or accumulate large steering bias. A
practical system must therefore regulate when and how strongly visual evidence
may modify the global command. It should also manage mucus during inspection
without discarding route progress.

We propose a progress-aware three-layer architecture (Fig.~\ref{fig:overview}).
Its contributions are:

\begin{enumerate}
  \item a compact goal-conditioned GRU trained from synchronized expert motor,
  visual, depth, and anatomical data, with directly auditable training metrics;
  \item a multimodal perception stack combining U-Net target segmentation,
  YOLO anatomical recognition, topology- and motion-constrained temporal
  confirmation, and Depth Anything V2 metric depth;
  \item a progress--distance weighted junction controller that inserts a
  bounded $M_1/M_2$ correction into the global command while preserving the
  insertion axis and long-horizon route direction;
  \item a continuously active suction strategy coupled with mucus tracking,
  depth-aware approach, saved navigation state, mission continuation, and
  autonomous homing; and
  \item an evaluation framework spanning model accuracy, route completion,
  centerline tracking, wall contact, mucus clearance, continuity, and optional
  bounded reinforcement-learning refinement.
\end{enumerate}

\begin{figure}[t]
\centering
\fbox{\parbox[c][4.3cm][c]{0.94\columnwidth}{\centering
\textbf{FIGURE 1 PLACEHOLDER --- Complete three-layer architecture}\\[1.5mm]
Draw five horizontal blocks: expert data and GRU training; RGB/YOLO/U-Net/Depth
Anything perception; Layer~1 global policy; Layer~2 progress-weighted visual
correction; Layer~3 mucus intervention and mission continuity. Show camera,
$M_0$--$M_2$, continuously active vacuum pump, motor feedback, anatomical goal,
saved progress index, and final homing arrow. Use different colors for learned
models, deterministic safety logic, and physical hardware.}}
\caption{The proposed multimodal AutoNav3Layer framework.}
\label{fig:overview}
\end{figure}

\section{Related Work}

\subsection{Robotic Bronchoscopy and Continuum Control}

Bronchoscope mechanics inherit the nonlinearities of tendon-driven continuum
robots. Constant-curvature, tendon-routing, Cosserat-rod, and medical continuum
robot studies provide the modeling foundation
\cite{camarillo2008mechanics,webster2010continuum,rucker2011statics,
burgner2015survey,till2019cosserat}. Friction, hysteresis, mounting variation,
and wall interaction limit model-only control. Deep visual servoing has thus
been explored for intraluminal navigation~\cite{lazo2022visualservo};
bronchoscopy-specific systems additionally use CT-derived environments, images,
or robot pose~\cite{zhang2024aicopilot,zhao2024bronchocopilot,xu2025lungdepth}.
Our approach separates global anatomical intent from bounded local correction.

\subsection{Imitation and Reinforcement Learning}

Learning from demonstration transfers expert skill without unsafe unrestricted
exploration~\cite{argall2009survey,osa2018imitation}. Behavioral cloning is
data efficient but can suffer from distribution shift, motivating aggregation
and adversarial formulations~\cite{ross2011dagger,ho2016gail}. Goal tokens
disambiguate actions at shared observations~\cite{codevilla2018conditional},
while LSTM and GRU models capture temporal context
\cite{hochreiter1997lstm,cho2014gru}. Recent robot-learning studies emphasize
demonstration quality, temporal action representations, multimodality, and
closed-loop evaluation
\cite{mandlekar2022robomimic,florence2022ibc,shafiullah2022bet,
chi2024diffusionpolicy,zhao2023act,lynch2020play,levine2016visuomotor}.
For conservative post-training refinement, we formulate a bounded PPO residual
using the clipped policy objective~\cite{schulman2017ppo} and a reproducible
implementation framework~\cite{raffin2021sb3}.

\subsection{Endoscopic Perception}

U-Net and later segmentation architectures established strong dense-prediction
baselines~\cite{ronneberger2015unet,he2017maskrcnn,chen2018deeplab}. Foundation
models improve transferable visual representations
\cite{kirillov2023sam,oquab2024dinov2}, while real-time one-stage detection
supports anatomical recognition~\cite{redmon2016yolo}. Monocular depth and
endoscopic adaptation provide geometric cues under limited texture
\cite{godard2019monodepth2,yang2024depthanythingv2,ozyoruk2021endoslam,
cui2024endodac,zeinoddin2024dares}. Because airway suction must be bounded and
clinically supervised, we also follow bronchoscope suction studies and airway
management guidance~\cite{lamb2022suction,blakeman2022suction}.

\section{System and Data}

\subsection{Platform}

The system comprises a flexible bronchoscope, RGB endoscopic camera, three
position-controlled motors, encoders, a vacuum suction pump, and a GPU
workstation. $M_0$ controls insertion; $M_1$ and $M_2$ bend the distal tip in
orthogonal directions. The pump remains active throughout autonomous inspection,
so Layer~3 controls alignment and exposure rather than pump switching. Position,
speed, image freshness, depth validity, and operator-stop conditions are checked
before every command.

\begin{table}[t]
\caption{Hardware and Runtime Configuration}
\label{tab:hardware}
\centering\small
\resizebox{\columnwidth}{!}{%
\begin{tabular}{lll}
\toprule
Component & Role & Final reported value \\
\midrule
$M_0$ & insertion & range $[-900,0]^\circ$, max $80^\circ$/s \\
$M_1$ & horizontal bending & range $[-170,170]^\circ$, max $30^\circ$/s \\
$M_2$ & vertical bending & range $[-500,500]^\circ$, max $80^\circ$/s \\
RGB camera & endoscopic perception & \RESULT{resolution, measured fps} \\
Vacuum pump & continuous secretion suction & \RESULT{model, pressure/flow} \\
GPU/CPU & model inference & \RESULT{hardware and total latency} \\
Phantom & airway evaluation & \RESULT{model/material/scale} \\
\bottomrule
\end{tabular}}
\end{table}

\subsection{Expert Demonstrations}

Manual expert operation was synchronized with images, motor angles and
velocities, perception features, depth values, timestamps, anatomical targets,
and motor increments. The curated set contains 87 complete trajectories and
30,730 frames over seven recorded targets (Table~\ref{tab:data}). Splitting is
performed by whole trajectory to prevent adjacent-frame leakage.

\begin{table}[t]
\caption{Goal-Conditioned Expert Dataset}
\label{tab:data}
\centering\small
\resizebox{\columnwidth}{!}{%
\begin{tabular}{lrrr}
\toprule
Target & Trajectories & Frames & Mean frames \\
\midrule
TR & 5 & 640 & 128 \\
LMB & 23 & 8,196 & 356 \\
RMB & 26 & 6,326 & 243 \\
LUB & 8 & 3,344 & 418 \\
LLB & 8 & 6,071 & 759 \\
RML & 8 & 3,237 & 405 \\
RLL & 9 & 2,916 & 324 \\
\midrule
Total & 87 & 30,730 & 353 \\
\bottomrule
\end{tabular}}
\end{table}

\section{Multimodal Perception}

\subsection{U-Net Mucus Segmentation}

The deployed segmentation branch uses a VGG-encoder U-Net with skip-connected
decoder blocks and a $1\!\times\!1$ output projection. RGB frames are resized to
$512\!\times\!512$. The current inference configuration separates mucus from
background; a three-class training variant additionally represents the
instrument. Training uses cross-entropy plus soft Dice loss,
\begin{equation}
 \mathcal{L}_{seg}=\mathcal{L}_{CE}+\lambda_D(1-\mathrm{Dice}),
\end{equation}
Adam with initial learning rate $10^{-4}$, cosine decay, and batch size two.
Geometric and photometric augmentation should be reported exactly from the
training run. The largest valid connected mucus component supplies centroid
$\mathbf{c}^m_t$, area $A^m_t$, and mask $M^m_t$.

\begin{figure}[t]
\centering
\fbox{\parbox[c][4.0cm][c]{0.94\columnwidth}{\centering
\textbf{FIGURE 2 PLACEHOLDER --- U-Net training and segmentation}\\[1.5mm]
Panel (a): representative RGB/mask pairs for clean lumen, thin mucus, thick
mucus, specular highlight, blur, and partial occlusion. Panel (b): VGG-U-Net
encoder/decoder diagram. Panel (c): train/validation loss and Dice curves.
Panel (d): test examples with TP/FP/FN overlays. Include dataset sizes and 95\%
CI for Dice and IoU.}}
\caption{Mucus segmentation training and held-out performance.}
\label{fig:unet}
\end{figure}

\subsection{YOLO Anatomical Recognition and Temporal Confirmation}

A lightweight nine-class YOLO detector recognizes TR, RMB, RUL, RML, RLL, LMB,
LUB, LLB, and BI. The model was initialized from YOLO11n and trained for 120
epochs at $640\!\times\!640$, batch size 16, AdamW learning rate $5\times10^{-4}$,
cosine decay, and weak geometric/color augmentation. It contains 2.59M
parameters. Raw detections are filtered at confidence 0.30 and passed to a
motion-aware anatomical state estimator. Only the current node, an allowed
child during insertion, or the parent during withdrawal can be confirmed.
Candidates require repeated observations; evidence is frozen during stationary
frames, and branch-specific windows use insertion progress and bending evidence.
This prevents single-frame anatomical jumps and makes location estimates
consistent with the bronchial tree.

\begin{table}[t]
\caption{Verified YOLO Validation Metrics}
\label{tab:yolo}
\centering\small
\begin{tabular}{lcccc}
\toprule
Classes & Precision & Recall & mAP50 & mAP50--95 \\
\midrule
9 & 0.8926 & 0.7797 & 0.9096 & 0.6006 \\
\bottomrule
\end{tabular}
\end{table}

\subsection{Depth Anything V2 and Robust Target Depth}

Metric Depth Anything V2 estimates a dense depth map $D_t$. For a target mask,
we erode boundary pixels, reject invalid samples, compute a robust central depth,
apply scale/bias calibration, reject large temporal jumps, and filter with an
EMA coefficient of 0.2. The resulting depth $\bar z_t$ controls insertion:
approach when the target is distant, stop inside a safety band, and withdraw
when too close. Depth is used as a control cue rather than an uncalibrated
absolute measurement; phantom calibration error must be reported.

\begin{figure}[t]
\centering
\fbox{\parbox[c][4.1cm][c]{0.94\columnwidth}{\centering
\textbf{FIGURE 3 PLACEHOLDER --- Multimodal perception results}\\[1.5mm]
Use four columns for raw RGB, U-Net mucus mask, YOLO anatomical detection, and
Depth Anything map. Add a topology panel showing raw YOLO label, candidate
counter, confirmed label, insertion/withdrawal direction, and permitted tree
transition. Add a depth-calibration scatter plot with fitted line and error.}}
\caption{U-Net, YOLO, topology, and depth outputs in representative airway scenes.}
\label{fig:perception}
\end{figure}

\section{Hierarchical Control}

\subsection{Layer 1: Goal-Conditioned Recurrent Imitation}

Each observation $\obs_t\in\mathbb{R}^{20}$ contains mucus and junction
features, center/mean/maximum depth, three directional path depths, three motor
angles, three motor velocities, and elapsed time. An MLP maps it through
$20\!\rightarrow\!128\!\rightarrow\!128$ with LayerNorm and ELU. The anatomical
goal is embedded in 16 dimensions and concatenated across an eight-frame window:
\begin{align}
 \mathbf{e}_i&=f_o(\obs_i),\quad
 \mathbf{u}_i=[\mathbf{e}_i\|\mathrm{Embed}(\goal)],\\
 \hid_i&=\mathrm{GRU}(\mathbf{u}_i,\hid_{i-1}),\quad i=t-7,\ldots,t.
\end{align}
A two-layer GRU has hidden size 256 and dropout 0.1. Its action head is
$256\!\rightarrow\!128\!\rightarrow\!3$ with ELU and $\tanh$; the task head is
$256\!\rightarrow\!64\!\rightarrow\!2$. Normalized actions are scaled by
$[5,3,3]^\circ$ for $M_0/M_1/M_2$. The total training objective is
\begin{equation}
 \mathcal{L}_{BC}=\|\hat\act_t-\act^E_t\|_2^2+
 \lambda_s\mathrm{CE}(\hat{\mathbf z}_t,s^E_t).
\end{equation}
Training uses AdamW, OneCycleLR, mixed precision, gradient norm clipping at 1.0,
batch size 64, and a trajectory-wise 90/10 split. The best model is selected by
validation navigation loss.

\begin{figure}[t]
\centering
\fbox{\parbox[c][4.1cm][c]{0.94\columnwidth}{\centering
\textbf{FIGURE 4 PLACEHOLDER --- GRU structure and training curves}\\[1.5mm]
Left: exact tensor dimensions for the 20-D observation encoder, 16-D goal
embedding, 8-frame GRU, action head, and task head. Right: train/validation
navigation loss and task accuracy for all available runs; mark the selected
epoch 29 model and report mean$\pm$SD over at least three seeds.}}
\caption{Goal-conditioned recurrent imitation policy and optimization history.}
\label{fig:gru}
\end{figure}

\begin{table}[t]
\caption{Verified GRU Training Evidence}
\label{tab:gru}
\centering\small
\resizebox{\columnwidth}{!}{%
\begin{tabular}{lcccc}
\toprule
Model & Parameters & Selected epoch & Val. action MSE & Val. task acc. \\
\midrule
Goal-conditioned GRU & 773,221 & 29 & 0.002084 & 1.000 \\
\bottomrule
\end{tabular}}
\end{table}

Layer~1 integrates the action into a bounded nominal target. To retain an
auditable global direction under distribution shift, the controller constrains
the learned command inside an expert-consistent envelope. The exact deployment
coefficient and envelope width must be logged in the final experiment. This
separation allows the learned policy to express goal-conditioned motor timing
while preserving safe long-horizon intent.

\subsection{Layer 2: Progress-Weighted Junction Correction}

Dark-lumen candidates are extracted after circular masking and morphological
opening. Among candidates exceeding the minimum area, the geometric center
nearest the image center is selected. For normalized error
$\mathbf e_t=[e^x_t,e^y_t]^\top$, the correction authority is the product of a
progress weight and distance weight. With route progress $s_t\in[0,1]$,
\begin{equation}
 w_p(s_t)=\beta_0+(1-\beta_0)
 \min\!\left(\frac{s_t}{s_f},1\right)^2,
\end{equation}
where $\beta_0=0.25$ and $s_f=0.80$. For center distance
$d_t=\|\mathbf e_t\|_2$,
\begin{equation}
 w_d(d_t)=\begin{cases}
 1,&d_t\le0.30,\\
 (0.75-d_t)/(0.75-0.30),&0.30<d_t<0.75,\\
 0,&d_t\ge0.75.
 \end{cases}
\end{equation}
After an axis dead-zone of 0.08, raw offsets use gains
$K_1=3.0$ and $K_2=2.5$, EMA coefficient 0.25, and per-axis total limit
$5^\circ$. The command is
\begin{equation}
 \qcmd_t=\qnom_t+
 w_p(s_t)w_d(d_t)[0,\delta q^v_{1,t},\delta q^v_{2,t}]^\top.
\end{equation}
Thus $M_0$ preserves Layer~1 progress, whereas $M_1/M_2$ receive small,
smooth, nonaccumulating centerline corrections. Independent visual alignment
uses a 10-pixel dead-zone and a far-fast/near-slow speed curve; motor-axis signs
and vertical scaling were calibrated on the physical robot. This structure is
related to bounded residual control~\cite{johannink2019residual}, but uniquely
modulates authority with procedural progress and visual geometry.

\begin{figure}[t]
\centering
\fbox{\parbox[c][4.1cm][c]{0.94\columnwidth}{\centering
\textbf{FIGURE 5 PLACEHOLDER --- Layer 2 correction mechanism}\\[1.5mm]
Panel (a): multiple lumen candidates and nearest-center selection. Panel (b):
$w_p(s)$ and $w_d(d)$ curves. Panel (c): Layer~1 nominal $M_1/M_2$, bounded
visual offsets, and final commands across a bifurcation. Panel (d): image-plane
center error with/without Layer~2 and synchronized wall-contact markers.}}
\caption{Continuous progress--distance weighting preserves global intent while improving local centering.}
\label{fig:layer2}
\end{figure}

\subsection{Layer 3: Continuous Suction and Mission Continuity}

The vacuum pump remains active during the inspection. When the U-Net mask
exceeds the valid-area threshold, the controller stores route progress $k^*$
and motor configuration $\mathbf q^*$, then prioritizes mucus tracking.
$M_1/M_2$ center the mask, while $M_0$ approaches according to filtered depth
and stops within the safety band. Missing detections do not immediately end the
intervention: the target must remain absent for $N_{clear}=15$ frames. This
dropout-tolerant rule, together with mask-area evolution and depth validity,
suppresses false completion caused by blur or specular highlights.

After clearance, the robot returns to $\mathbf q^*$ within a $5^\circ$ maximum
axis error, re-enters the remaining route at progress index $k^*$, completes the
selected airway inspection, and finally drives all three navigation axes to the
defined origin. The hierarchy is therefore
\begin{align}
 \textsc{Navigate}&\rightarrow\textsc{Track/Clear}\rightarrow\textsc{Saved Pose}\nonumber\\
 &\rightarrow\textsc{Continue}\rightarrow\textsc{Home}.
\end{align}
Continuous suction simplifies actuator coordination but requires reporting
pump pressure, total exposure time, and any mucosal-contact surrogate.

\begin{figure}[t]
\centering
\fbox{\parbox[c][4.3cm][c]{0.94\columnwidth}{\centering
\textbf{FIGURE 6 PLACEHOLDER --- Layer 3 sequence and state variables}\\[1.5mm]
Show six real frames: cruising with pump on, mucus detection, depth-aware
approach, centered suction, 15-frame absence confirmation, and continued route.
Below, plot route index, $M_0$--$M_2$, mask area, depth, target error, pump state,
and final homing. Mark the saved pose and maximum return error.}}
\caption{Continuous suction, mucus tracking, progress-aware mission continuity, and homing.}
\label{fig:layer3}
\end{figure}

\subsection{Bounded Reinforcement-Learning Refinement}

We formulate PPO as an optional residual tuner rather than an unrestricted
motor policy. Its state contains GRU hidden features, route progress, center
error, depth, tracking error, and recent contact evidence. Its action adjusts
only Layer~2 gain and offset within fixed bounds. A suitable reward is
\begin{align}
 r_t={}&w_g\Delta s_t-w_c C_t-w_e\|\mathbf e_t\|_2
 -w_j\|\Delta\act_t\|_2^2\nonumber\\
 &-w_b\|\delta\mathbf q^v_t\|_2^2+w_T\mathbb I_{goal},
\end{align}
where $C_t$ is contact. Domain randomization covers start pose, illumination,
friction, mask noise, and depth scale. The learned residual is accepted only
when it improves held-out simulated success and then passes the same hardware
limits. The existing standalone PPO artifact is not sufficient for this claim:
its logged evaluation success is 0\%, so the AutoNav-specific PPO row remains a
required experiment rather than a reported improvement.

\section{Experiments and Results}

\subsection{Study Design}

Experiments should use paired, randomized trials on the same airway phantom.
Targets include at least LMB, RMB, LUB, LLB, RUL, RML, and RLL. Conditions are
nominal start, start-pose perturbation, illumination variation, lens
contamination, and mucus placement at predefined locations. Report sample size
per target, failure taxonomy, 95\% confidence intervals, and paired tests with
multiple-comparison correction.

Baselines are: (B0) expert manual control; (B1) frame-wise BC; (B2) GRU without
goal; (B3) goal-conditioned GRU/Layer~1; (B4) reactive visual tracking only;
(B5) Layer~1+2; (B6) full Layer~1+2+3; and (B7) full system with validated PPO
residual, if available. Metrics are route success, wrong-branch rate, centerline
error, contact count/duration, completion time, tracking error, intervention
count, mucus clearance, clearance time, return-to-saved-pose error, continued
route success, and homing error.

\begin{figure}[t]
\centering
\fbox{\parbox[c][3.8cm][c]{0.94\columnwidth}{\centering
\textbf{FIGURE 7 PLACEHOLDER --- Phantom protocol}\\[1.5mm]
Photograph the robot, fixed start jig, airway phantom, target branches, vacuum
pump and collection bottle. Overlay predefined mucus positions and perturbation
conditions. Add a randomized trial timeline and failure taxonomy.}}
\caption{Standardized hardware evaluation setup.}
\label{fig:setup}
\end{figure}

\subsection{Model-Level Results}

Table~\ref{tab:gru} and Table~\ref{tab:yolo} contain values recovered directly
from saved model artifacts. The selected GRU model reduced validation action
MSE from 0.02577 in the first logged epoch to 0.002084 at epoch 29, while task
accuracy remained 100\% on the recorded validation split. Because task labels
may be easy to separate, macro-F1 and per-class confusion must also be reported.
The YOLO detector reached mAP50 0.9096, but recall 0.7797 indicates that temporal
confirmation must tolerate missed frames. U-Net and depth results cannot be
completed until the held-out masks and calibration logs are restored.

\begin{table}[t]
\caption{Perception Model Results to Report}
\label{tab:perception}
\centering\small
\resizebox{\columnwidth}{!}{%
\begin{tabular}{lccccc}
\toprule
Model & Task & Precision & Recall & Dice/mAP & Latency \\
\midrule
VGG-U-Net & mucus mask & \RESULT{--} & \RESULT{--} & \RESULT{Dice/IoU} & \RESULT{ms} \\
YOLO11n & anatomy & 0.8926 & 0.7797 & 0.9096 mAP50 & \RESULT{ms} \\
Depth Anything V2 & depth & -- & -- & \RESULT{MAE/RMSE} & \RESULT{ms} \\
Fusion stack & synchronized & -- & -- & \RESULT{valid-frame rate} & \RESULT{fps} \\
\bottomrule
\end{tabular}}
\end{table}

\subsection{Integrated Results}

Qualitative integrated tests completed the full sequence of target-directed
navigation, progress-weighted junction correction, continuous mucus suction,
mission continuation, route completion, and homing. The final paper must replace
this statement with the quantitative trial matrix below; videos alone are not
sufficient evidence.

\begin{table}[t]
\caption{Navigation and Junction-Centering Comparison}
\label{tab:navigation}
\centering\small
\resizebox{\columnwidth}{!}{%
\begin{tabular}{lccccc}
\toprule
Method & Success & Wrong branch & Center error & Contacts & Time \\
\midrule
Expert & \RESULT{--} & \RESULT{--} & \RESULT{--} & \RESULT{--} & \RESULT{--} \\
Frame BC & \RESULT{--} & \RESULT{--} & \RESULT{--} & \RESULT{--} & \RESULT{--} \\
Layer~1 & \RESULT{--} & \RESULT{--} & \RESULT{--} & \RESULT{--} & \RESULT{--} \\
Reactive vision & \RESULT{--} & \RESULT{--} & \RESULT{--} & \RESULT{--} & \RESULT{--} \\
Layer~1+2 & \RESULT{--} & \RESULT{--} & \RESULT{--} & \RESULT{--} & \RESULT{--} \\
\bottomrule
\end{tabular}}
\end{table}

\begin{table}[t]
\caption{Ablation of Layer~2 Optimization}
\label{tab:layer2}
\centering\small
\resizebox{\columnwidth}{!}{%
\begin{tabular}{lcccc}
\toprule
Variant & Center error & Contacts & Max offset & Success \\
\midrule
Largest lumen & \RESULT{--} & \RESULT{--} & \RESULT{--} & \RESULT{--} \\
Nearest lumen & \RESULT{--} & \RESULT{--} & \RESULT{--} & \RESULT{--} \\
+ dead-zone/EMA & \RESULT{--} & \RESULT{--} & \RESULT{--} & \RESULT{--} \\
+ progress weight & \RESULT{--} & \RESULT{--} & \RESULT{--} & \RESULT{--} \\
Full bounded method & \RESULT{--} & \RESULT{--} & $\le5^\circ$ & \RESULT{--} \\
\bottomrule
\end{tabular}}
\end{table}

\begin{table}[t]
\caption{Mucus Intervention and Mission Continuity}
\label{tab:layer3}
\centering\small
\resizebox{\columnwidth}{!}{%
\begin{tabular}{lccccc}
\toprule
Method & Clear success & Clear time & Pose error & Continue success & Home error \\
\midrule
Manual suction & \RESULT{--} & \RESULT{--} & -- & \RESULT{--} & \RESULT{--} \\
Tracking, no depth & \RESULT{--} & \RESULT{--} & \RESULT{--} & \RESULT{--} & \RESULT{--} \\
Tracking + depth & \RESULT{--} & \RESULT{--} & \RESULT{--} & \RESULT{--} & \RESULT{--} \\
Full Layer~3 & \RESULT{--} & \RESULT{--} & \RESULT{--} & \RESULT{--} & \RESULT{--} \\
\bottomrule
\end{tabular}}
\end{table}

\begin{table}[t]
\caption{Reinforcement-Learning Refinement}
\label{tab:rl}
\centering\small
\resizebox{\columnwidth}{!}{%
\begin{tabular}{lccc}
\toprule
Policy & Sim. success & Hardware success & Contact reduction \\
\midrule
Standalone prototype & 0\% & not tested & not established \\
Bounded AutoNav PPO & \RESULT{--} & \RESULT{--} & \RESULT{--} \\
\bottomrule
\end{tabular}
}
\end{table}

\begin{figure}[t]
\centering
\fbox{\parbox[c][4.3cm][c]{0.94\columnwidth}{\centering
\textbf{FIGURE 8 PLACEHOLDER --- Final quantitative results}\\[1.5mm]
Panel (a): route success by target and method with 95\% CI. Panel (b): paired
center error and contact duration for Layer~1 vs Layer~1+2. Panel (c): mucus
area/depth/centering through Layer~3. Panel (d): saved-pose error, continued
route success and homing error. Panel (e): ablation waterfall.}}
\caption{Required final system-level results and ablations.}
\label{fig:results}
\end{figure}

\subsection{Software Verification}

The earlier control revision passed 92 hardware-independent logical checks.
The current anatomy-recognition interface additionally passes 20 logic tests
covering motor-direction gating, stale feedback, paused-state evidence,
anatomical topology, latest-frame transport, detection/recording separation,
and navigation-menu behavior. Layer~2 also has dedicated finite-offset and
long-duration tests; one source-consistency assertion currently expects an
older vertical speed scale (4.0 rather than the tuned 2.0) and must be updated
before the paper claims a fully clean regression suite. These checks are not
physical trials and are excluded from efficacy statistics.

\section{Discussion}

The system's main novelty is not any single network, but authority-aware fusion.
The GRU captures goal-conditioned temporal motor behavior; YOLO maintains
anatomical context; U-Net and depth supply local therapeutic geometry; and
Layer~2 inserts only a bounded correction whose strength depends on route
progress and visual distance. This avoids two common failure modes: a local
controller overriding the intended branch and a global controller ignoring
current wall geometry.

The current model evidence also reveals limitations. GRU task accuracy of 100\%
may reflect an easy label distribution and needs macro-F1 and target-stratified
testing. YOLO recall below 0.8 justifies temporal confirmation but may delay
anatomical transitions. U-Net metrics and depth calibration are missing from
the repository. Continuous suction reduces switching complexity but requires
pressure, exposure, fluid-removal, and tissue-safety reporting. Finally, the
standalone PPO run did not solve its evaluation task; reinforcement learning
must not be credited with system improvement until the bounded AutoNav residual
passes both simulation and hardware ablations.

Future work will expand expert data to missing/rare goals, collect multi-operator
demonstrations, train U-Net with hard negatives and temporal consistency,
calibrate depth across airway materials, learn uncertainty-aware fusion, and
evaluate on ex-vivo and in-vivo models. Clinical translation will additionally
require sterility, suction-pressure safety, failure containment, and supervised
human takeover protocols.

\section{Conclusion}

We introduced a progress-aware three-layer framework for autonomous robotic
bronchoscopy. A goal-conditioned GRU learns temporal expert control; multimodal
perception combines U-Net mucus masks, YOLO anatomical recognition, and Depth
Anything geometry; a bounded progress-weighted visual controller improves
junction centering; and continuous suction is integrated with mucus tracking,
mission continuity, and homing. Saved artifacts provide quantitative GRU and
YOLO evidence, while integrated trials establish qualitative feasibility of the
complete procedure. The remaining tables explicitly define the quantitative
evidence required for an IEEE RA-L submission.

\bibliographystyle{IEEEtran}
\bibliography{literature/refs_verified}
\end{document}
